\section{Results}
\label{sec:results}

This section reports the hard-smoke evaluation on verified PyBugHive \texttt{black} debugging tasks. We treat the first six tasks as the primary quantitative result because they were run under the same model and harness configuration. The seventh verified task, \texttt{pybughive\_black\_234}, was completed later with \texttt{gpt-5.5} after the original \texttt{gpt-5.3-codex} configuration became unavailable, so it is reported only as exploratory supplementary evidence.

\subsection{Main same-model result}
\label{sec:results_main}

Table~\ref{tab:hard_first6_main} summarizes the same-model hard-smoke result over six held-out PyBugHive \texttt{black} tasks. Auto SKILL.md solved 5/6 tasks, matching the best solve rate achieved by the generic checklist and Reflexion-style memory. This is not evidence of solve-rate dominance: the strongest three methods are tied on solved tasks. The result is instead an efficiency result. Among solved tasks, Auto SKILL.md used 56,867 tokens per solved task, compared with 65,915 for Reflexion memory, 71,694 for the generic checklist, 75,231 for length-matched Reflexion, and 92,031 for the no-memory baseline.

\FloatBarrier
\begin{table}[t]
\centering
\small
\caption{Same-model hard-smoke results over the first six PyBugHive Black tasks.}
\label{tab:hard_first6_main}
\resizebox{\linewidth}{!}{%
\begin{tabular}{lrrrrrr}
\toprule
Method & Solved & Solve rate & Cycles/solved & Tokens/solved & Failed patches/run & Negative transfer \\
\midrule
Auto SKILL.md & 5/6 & 0.8333 & 2.00 & 56,867 & 0.83 & 0.1667 \\
Generic checklist & 5/6 & 0.8333 & 2.00 & 71,694 & 1.33 & 0.0000 \\
Reflexion memory & 5/6 & 0.8333 & 2.60 & 65,915 & 1.33 & 0.0000 \\
Length-matched Reflexion & 4/6 & 0.6667 & 2.75 & 75,231 & 1.50 & 0.0000 \\
No memory & 4/6 & 0.6667 & 3.00 & 92,031 & 1.17 & 0.0000 \\
\bottomrule
\end{tabular}
}
\end{table}


\begin{figure}[t]
\centering
% Requires \usepackage{svg}; alternatively convert the SVG to PDF and replace \includesvg with \includegraphics.
\includesvg[width=\linewidth]{figures/hard_first6_main_bars}
\Description{A grouped bar chart comparing five methods. Auto SKILL.md,
generic checklist, and Reflexion each solve five of six tasks, while Auto
SKILL.md has the lowest tokens per solved task.}
\caption{Same-model hard-smoke aggregate over the first six PyBugHive Black tasks. Auto SKILL.md matches the strongest solve rate while using fewer tokens per solved task than the memory and no-memory controls.}
\label{fig:hard_first6_main_bars}
\end{figure}

The comparison against the length-matched Reflexion control is especially important. Auto SKILL.md improves solve rate from 4/6 to 5/6 and reduces tokens per solved task by 24.4\%. Because the length-matched control supplies comparable additional text without the trigger/procedure/failure-mode structure, this result supports the interpretation that procedural organization matters beyond simply adding unstructured reflection text to the prompt. A stricter structure-versus-content control is reported below; it leads to a more cautious conclusion.

\FloatBarrier
\subsection{Process efficiency}
\label{sec:results_efficiency}

The process metrics show that Auto SKILL.md compresses successful debugging behavior. Relative to Reflexion memory, Auto SKILL.md reduces diagnosis-edit-test cycles from 2.60 to 2.00 while also reducing tokens per solved task. Relative to no memory, it improves solve rate from 4/6 to 5/6 and reduces token cost by 38.2\%. Relative to the generic checklist, it has the same solve rate and cycle average but uses 20.7\% fewer tokens.

The paired common-solved analysis makes this more conservative by comparing Auto SKILL.md against each baseline only on tasks both methods solved. On the four tasks solved by both Auto SKILL.md and Reflexion memory, Auto SKILL.md reduces cycles by 36.4\% and tokens by 32.2\%. Against length-matched Reflexion on the same common-solved set, Auto SKILL.md also reduces cycles by 36.4\% and tokens by 36.2\%. These paired comparisons are useful because solve-conditional token averages can otherwise be affected by different methods solving different subsets of tasks.

\begin{figure}[t]
\centering
\includesvg[width=\linewidth]{figures/hard_first6_efficiency_scatter}
\Description{A scatter plot of cycles per solved task against tokens per
solved task for five methods. Auto SKILL.md lies toward the lower-left,
indicating comparatively efficient successful repairs.}
\caption{Process efficiency among solved hard-smoke tasks. Points show each method's average diagnosis-edit-test cycles and tokens per solved task over the same-model first-six aggregate. Lower-left indicates more efficient successful repair behavior; marker size encodes solve rate.}
\label{fig:hard_first6_efficiency}
\end{figure}

\FloatBarrier
\subsection{Negative transfer and subfamily sensitivity}
\label{sec:results_negative_transfer}

Auto SKILL.md produced the only explicit negative-transfer flag in the same-model hard-smoke set. On \texttt{pybughive\_black\_132}, it failed and was annotated as performance negative transfer because the induced procedure directed attention toward a nearby but incorrect formatter invariant; the stored trajectory notes that the first attempted fix increased full-suite failures from four to six before later narrowing. This matters for the paper's framing: SKILL.md should not be presented as universally beneficial memory. A more defensible claim is that procedural memory can improve within-family transfer efficiency while still requiring explicit negative-transfer accounting.

\begin{figure}[t]
\centering
\includesvg[width=\linewidth]{figures/hard_first6_task_outcomes_heatmap}
\Description{A six-task by five-method outcome matrix. Most cells are solved;
Auto SKILL.md uniquely solves task black 193 and has a negative-transfer-marked
failure on task black 132.}
\caption{Per-task outcome matrix for the first-six hard-smoke tasks. Green cells indicate solved tasks, red cells indicate unsolved tasks, and NT marks the negative-transfer flag.}
\label{fig:hard_first6_task_outcomes}
\end{figure}

The primary task-level matrix shows this subfamily sensitivity. Auto SKILL.md is the only method that solves task 193 and is the most efficient on task 232, but it fails on task 132. Reflexion memory and the generic checklist solve task 132 but fail on task 193. Different memory formats therefore transfer different prior regularities, motivating measurement of both successful and negative transfer.

\FloatBarrier
\subsection{Supplementary structural-control rerun}
\label{sec:results_structural_control}

After the main run, we added a same-model \texttt{gpt-5.5} comparison between Auto SKILL.md and a format-shuffled SKILL.md on the same six tasks. This rerun does not replace the main five-method table, which used \texttt{gpt-5.3-codex}. It tests whether procedural organization still separates behavior when skill content is preserved but the three-section organization is disrupted.

\FloatBarrier
\begin{table}[t]
\centering
\caption{Same-model \texttt{gpt-5.5} structural-control rerun on the first-six hard-smoke tasks. Both methods solve all tasks; Auto SKILL.md uses fewer recorded cycles and failed patches, while the shuffled control uses fewer tokens in aggregate.}
\label{tab:hard_first6_structural_control_gpt55}
\resizebox{\linewidth}{!}{%
\begin{tabular}{lrrrrrr}
\toprule
Method & Solved & Solve rate & Cycles / solved & Tokens / solved & Failed patches / run & NT rate \\
\midrule
Auto SKILL.md & 6/6 & 100\% & 1.67 & 111,978 & 0.50 & 0\% \\
Format-shuffled SKILL.md & 6/6 & 100\% & 2.17 & 102,006 & 0.83 & 0\% \\
\bottomrule
\end{tabular}
}

\end{table}

Both Auto SKILL.md and format-shuffled SKILL.md solved all six tasks under \texttt{gpt-5.5}. Auto SKILL.md used fewer recorded cycles per solved task (1.67 vs. 2.17) and fewer failed patches per run (0.50 vs. 0.83), but it was not more token-efficient in aggregate (111,978 vs. 102,006 tokens per solved task). The token difference is driven mainly by \texttt{pybughive\_black\_193}, where Auto SKILL.md solved but used 258,676 tokens, indicating substantial internal patch search despite being recorded as a single final successful cycle.

This supplementary result weakens any claim that SKILL.md structure is necessary for solvability under a stronger model. The safer interpretation is that procedural organization changes process behavior, especially recorded cycles and patch churn, but the current evidence does not show uniform token-efficiency superiority over shuffled skill content.

\FloatBarrier
\subsection{Qualitative case studies}
\label{sec:results_cases}

Two tasks are particularly useful for the narrative in the primary five-method run. \path{pybughive_black_193} is the clearest positive transfer case for Auto SKILL.md in that table: it is the only method that solves the task, while both Reflexion baselines pursue a nearby but insufficient prior pattern. This makes it a good case study for why a structured trigger/procedure/failure-mode artifact can transfer better than unstructured reflection snippets in the original hard-smoke setting.

\texttt{pybughive\_black\_232} is a process-efficiency case in the primary run. All methods solve it, but Auto SKILL.md solves it in 1 cycle and 36,035 tokens, while Reflexion and length-matched Reflexion each require 3 cycles and more than 68,000 tokens. This is the cleanest example of the paper's core efficiency claim against the main memory baselines: even when solve rate does not separate methods, procedural memory can reduce the amount of diagnosis and patch iteration needed to converge.

\texttt{pybughive\_black\_132} should be discussed as the counterexample from the primary run. It prevents overclaiming and gives the negative-transfer metric a concrete role. The failure suggests that the induced SKILL.md can overfit a nearby procedural invariant when the held-out task belongs to an adjacent but different formatter subfamily. Because the later \texttt{gpt-5.5} rerun solves this task, the case should be framed as model- and run-specific negative transfer rather than as a stable property of the task.

\subsection{Exploratory all-seven result}
\label{sec:results_exploratory}

The all-seven aggregate preserves the same broad pattern: Auto SKILL.md, Reflexion memory, and the generic checklist each solve 6/7 tasks, while length-matched Reflexion and no memory solve 5/7. Auto SKILL.md remains the most token-efficient successful method among the five primary methods overall. However, this table is model-mixed because \texttt{pybughive\_black\_234} was run with \texttt{gpt-5.5}, so it should be reported as supplementary evidence rather than as the main quantitative result.

\subsection{Result takeaway}
\label{sec:results_takeaway}

The main empirical takeaway is that induced natural-language SKILL.md is best framed as a procedural memory artifact that improves low-shot within-family debugging efficiency. The current evidence does not show universal solve-rate superiority. Instead, it shows that Auto SKILL.md can match the best solve rate on a hard same-model slice while reducing token cost and diagnosis-edit-test cycles relative to unstructured reflection memories and no-memory controls, with an explicit negative-transfer failure that should be retained rather than hidden. The supplementary shuffled-control rerun narrows the structural claim: organization appears to affect process behavior, but the current evidence does not prove that the canonical SKILL.md structure is necessary for solvability.
